%=================================================================
% Supplementary Materials for
% "Composition-Level Constraints in Alternative Gravity:
%  A Recursive Constraint-Transport Criterion"
%=================================================================
\documentclass[11pt]{article}

\usepackage[margin=1in]{geometry}
\usepackage[T1]{fontenc}
\usepackage[utf8]{inputenc}
\usepackage{lmodern}
\usepackage{microtype}
\usepackage{amsmath,amssymb,amsthm,mathtools}
\usepackage{booktabs,array,longtable}
\usepackage{enumitem}
\usepackage{hyperref}
\hypersetup{hidelinks}
\setlength{\parskip}{0.45em}
\setlength{\parindent}{0pt}
\setlist{itemsep=0.3em,topsep=0.4em}

\title{Supplementary Material: A Finite Second-Order Theory of Inverse Gravity\\
\large Clock-Orientation Carrier, Pair Resolution, and Registration Limits}
\author{Andy E. Williams\thanks{Correspondence: info@cc4ci.org}\\[4pt]
\small Caribbean Center for Collective Intelligence, St.~John's, Antigua and Barbuda\\
\small Nobeah Foundation, Nairobi, Kenya}
\date{}

%----------------------- notation -----------------------
\newcommand{\Pspace}{\ensuremath{\mathcal P}}
\newcommand{\ARspace}{\ensuremath{\mathcal A_R}}
\newcommand{\Graph}{\ensuremath{\Gamma}}
\newcommand{\Kgroup}{\ensuremath{K}}
\newcommand{\Uone}{\ensuremath{\mathrm U(1)}}
\newcommand{\SOthree}{\ensuremath{\mathrm{SO}(3)}}
\newcommand{\Id}{\ensuremath{\mathbf 1}}
\newcommand{\zinv}{\ensuremath{\mathsf z}}
\newcommand{\State}{\ensuremath{\mathcal Q}}
\newcommand{\Defect}{\ensuremath{\mathcal E}}
\newcommand{\Flow}{\ensuremath{\mathsf{Flow}}}
\newcommand{\Resp}{\ensuremath{\mathsf{Resp}}}
\newcommand{\Reg}{\ensuremath{\mathsf{Reg}}}
\newcommand{\Obs}{\ensuremath{\mathsf{Obs}}}
\newcommand{\Bridge}{\ensuremath{\mathfrak B}}
\newcommand{\Quot}{\ensuremath{\mathsf q}}
\newcommand{\Branch}{\ensuremath{\mathsf B}}
\newcommand{\Pair}{\ensuremath{\Pi}}
\newcommand{\PairRes}{\ensuremath{\mathsf{PairRes}}}
\newcommand{\Resolve}{\ensuremath{\mathsf{Resolve}}}
\newcommand{\Sched}{\ensuremath{\mathsf{Sched}}}
\newcommand{\Adm}{\operatorname{Adm}}
\newcommand{\supp}{\mathsf{supp}}
\newcommand{\dist}{\operatorname{dist}}
\newcommand{\grad}{\operatorname{grad}}
\newcommand{\argmin}{\operatorname*{arg\,min}}
\newcommand{\Phys}{\ensuremath{\mathsf{Phys}}}
\newcommand{\Emp}{\ensuremath{\mathsf{Emp}}}
\newcommand{\GC}{\ensuremath{\mathsf{GC}}}
\newcommand{\Status}{\ensuremath{\mathsf{Status}}}
\newcommand{\Register}{\operatorname{Register}}
\newcommand{\Recall}{\operatorname{Recall}}
\newcommand{\GOp}{\ensuremath{\mathsf G}}
\newcommand{\LOp}{\ensuremath{\mathsf L}}
\newcommand{\Inv}{\ensuremath{\mathsf{Inv}}}

% claim tags
\newcommand{\Defn}{\textsc{[def]}}
\newcommand{\Asm}{\textsc{[asm]}}
\newcommand{\Hyp}{\textsc{[hyp]}}
\newcommand{\Thm}{\textsc{[thm]}}
\newcommand{\Cnd}{\textsc{[cnd]}}
\newcommand{\Open}{\textsc{[open]}}

\theoremstyle{definition}
\newtheorem{definition}{Definition}[section]
\newtheorem{hypothesis}[definition]{Hypothesis}
\newtheorem{construction}[definition]{Construction}
\newtheorem{proposition}[definition]{Proposition}
\newtheorem{theorem}[definition]{Theorem}
\newtheorem{corollary}[definition]{Corollary}
\newtheorem{remark}[definition]{Remark}
\newtheorem{obligation}[definition]{Empirical obligation}

\begin{document}
\maketitle

\begin{abstract}
This supplement supplies a self-contained finite second-order theory of inverse
gravity.  Its first-order physical carrier is a finite matter-generated
clock-orientation gravity model with carrier
$K=\mathrm U(1)\times\mathrm{SO}(3)$.  A primitive counterposed interaction
pair has an intrinsic transport-relative separation and an intrinsic
synchronization threshold.  Below the threshold the pair is physically
co-registered and selects the regular attractive branch; above it the same
counterposed pair is physically resolved and selects an inverse-response branch.
One finite energy and its gradient flow then give the forced branch responses:
regular configurations have a stable attractive separation, while inverse
configurations satisfy a positive separation response.

The second-order layer is a distinct registration theory.  A typed bridge maps
P-space states into material records and a quotient retains only the record
available to a declared observer.  The resulting factorization theorem states
exactly when a material record can decide the physical branch.  Thus an inverse
response may be physically present while an admissible local record co-registers
it with an attractive state.  The supplement also states the role of a
target-conditioned schedule: it is a conditional realization or preparation
object for an already-defined inverse branch, not the definition of inverse
gravity.  The theory is complete on its declared finite model domain; physical
realization, continuum recovery, and empirical selection remain separately typed
questions.
\end{abstract}

\tableofcontents
\newpage

\section{Reading guide, status, and space discipline}
\label{sec:guide}

The main article studies a decision problem: when does a material observation
record preserve the relation needed to decide a physical target?  The present
supplement supplies the missing physical object for that problem.  The object has
two distinct spaces:
\begin{equation}
\Pspace \not\equiv \ARspace.
\label{eq:spaces-distinct}
\end{equation}
\Pspace{} contains physical carriers, interaction pairs, transports, branch
conditions, and responses.  $\ARspace$ contains registrations, retained records,
resolution budgets, and quotients.  A conclusion may move between them only
through an explicit bridge.

The theory uses two complete generator inventories without identifying them.
The physical interaction inventory in \Pspace{} is
\begin{equation}
\mathcal I_P=
\{\text{strong},\text{weak},\text{electromagnetic},\text{gravity}\}.
\label{eq:physical-inventory}
\end{equation}
The operational grammar in $\ARspace$ is
\begin{equation}
\mathcal O_{A(R)}=\{\Register,\Recall,\GOp,\LOp\}.
\label{eq:registration-grammar}
\end{equation}
The present construction uses $\GOp$ and $\LOp$ only as the expansion and
contraction component of the registration/topological grammar.  They are not
asserted to replace the four physical interactions in Equation~\eqref{eq:physical-inventory}.

Every numbered statement is tagged by logical form.  A \Defn{} fixes an object;
an \Asm{} is a constitutive condition of this declared model class; a \Hyp{} is a
physical hypothesis; a \Thm{} follows from the stated definitions and
assumptions; a \Cnd{} states a conditional result; and an \Open{} entry records a
separate realization or empirical obligation.

\begin{definition}[Status record]
\label{def:status}
\Defn{} For a theory object $H$, write
\begin{equation}
\Status(H)=\bigl(\mathsf{Form}(H),\ \mathsf{Domain}(H),\
\mathsf{Bridge}(H),\ \mathsf{Selection}(H)\bigr).
\label{eq:status}
\end{equation}
The first coordinate distinguishes definition, hypothesis, theorem, and
conditional consequence.  The second names the finite carrier and regime.  The
third records the P-space--$A(R)$-space bridges on which any observability claim
depends.  The fourth records whether the object is mathematically defined,
physically realized, or empirically selected.  A finite theory may be complete
on its declared domain while its physical realization and empirical selection
remain open.
\end{definition}

\section{The finite P-space carrier}
\label{sec:carrier}

\subsection{Clock-orientation data}

\begin{definition}[Admissible finite gravity datum]
\label{def:datum}
\Defn{} An admissible finite datum is
\begin{equation}
\mathbf p=\bigl(
\Graph,T,E^+,\mathcal C_T,K,m,\xi,p,\mathsf S,
\Pair,\eta,c,h,\alpha,\gamma,\lambda,\kappa,\Lambda
\bigr),
\label{eq:datum}
\end{equation}
with the following components.
\begin{enumerate}
\item $\Graph=(V,E)$ is a finite connected graph, $E^+$ selects one orientation
of each unoriented edge, $T\subseteq E$ is a spanning tree, and $\mathcal C_T$
is the finite family of rooted, ordered fundamental cycles induced by the chords
of $T$.
\item The clock-orientation carrier is
\begin{equation}
K:=\Uone\times\SOthree,
\qquad
\zinv:=(-1,I_3)\in K,
\qquad \zinv^2=\Id.
\label{eq:K}
\end{equation}
$K$ contains one clock phase and one spatial orientation.  It does not posit a
full spacetime-frame group.
\item Each node $v$ carries a matter closure
$\eta_v=(m_v,\xi_v)\in(0,\infty)\times X_v$, where $X_v$ is a declared
finite-dimensional internal carrier.  Each oriented edge $e=(u,v)$ carries a
finite port datum $p_e$ and a total, matter-sensitive source map
\begin{equation}
U_e=\mathsf S_e(\eta_u,\eta_v,p_e)\in K,
\qquad U_{\bar e}=U_e^{-1}.
\label{eq:source-map}
\end{equation}
Matter sensitivity means that on every gravitational edge the source map changes
for at least one admissible change of endpoint closure data.
\item Every edge carries a counterposed primitive interaction pair
\begin{equation}
\Pair_e=(a_e^+,a_e^-,\delta_e,\vartheta_e),
\qquad \mathsf{Opp}_e(a_e^+,a_e^-)=1.
\label{eq:pair}
\end{equation}
Here $\delta_e\ge0$ is the pair's declared transport-relative separation and
$\vartheta_e$ retains the ordered interaction history needed to evaluate that
separation.  Neither is a raw global coordinate.
\item Every edge has an intrinsic synchronization threshold $\eta_e>0$ and a
central branch target $c_e\in\{\Id,\zinv\}$ defined in
Hypothesis~\ref{hyp:pair-resolution}.  Every fundamental cycle has a fixed
sector value $h_C\in K$.  Finally,
$\alpha_e,\gamma_e,\lambda_e,\kappa_e>0$ and $\Lambda_C\ge0$ are declared
finite coefficients.
\end{enumerate}
\end{definition}

\begin{hypothesis}[Primitive counterposition and intrinsic pair resolution]
\label{hyp:pair-resolution}
\Hyp{} At the declared primitive scale, a gravitational interaction is represented
by a counterposed pair as in Equation~\eqref{eq:pair}.  Its physical branch is
not decided by an observer.  It is decided in \Pspace{} by the intrinsic
transport-relative pair-resolution map
\begin{equation}
\PairRes_e(\Pair_e)=c_e:=
\begin{cases}
\Id, & 0\le \delta_e\le \eta_e,\\[2mm]
\zinv, & \delta_e>\eta_e.
\end{cases}
\label{eq:pair-resolution}
\end{equation}
The case $c_e=\Id$ is the co-registered branch: the counterposed pair is
unresolved relative to the intrinsic synchronization threshold.  The case
$c_e=\zinv$ is the resolved branch: the counterposed pair remains counterposed
but its separation is physically retained by the carrier.
\end{hypothesis}

The word ``intrinsic'' in Hypothesis~\ref{hyp:pair-resolution} is decisive.  The
threshold $\eta_e$ belongs to the declared P-space carrier.  It is not a detector
resolution and does not change merely because a different material observer is
used.  Detector noise appears later, in $\ARspace$.

\begin{remark}[Meaning of ``simultaneous'']
\label{rem:simultaneous}
\Defn{} In this supplement, ``simultaneous'' means co-registered below the
intrinsic transport-relative threshold $\eta_e$.  It does not assert a
frame-independent global simultaneity relation.  The primitive input is the
ordered pair history $\vartheta_e$, together with its declared separation
functional $\delta_e$.
\end{remark}

\subsection{State, transport, and integrability}

\begin{definition}[State space and relative mismatch]
\label{def:state}
\Defn{} For an admissible datum $\mathbf p$, the total finite state space is
\begin{equation}
\State_{\mathbf p}:=K^V\times K^{E^+}\times(0,\infty)^{E^+}.
\label{eq:state-space}
\end{equation}
A state is $z=(q,W,r)$, with node clock orientations $q_v\in K$, actual
edge transports $W_e\in K$, and positive interaction separations $r_e>0$.
The transport-relative mismatch is
\begin{equation}
\mu_e(q,W):=q_v^{-1}W_eq_u\in K
\qquad(e=(u,v)\in E^+).
\label{eq:mismatch}
\end{equation}
A datum is source-target integrable in the identity sector when there is a
$q^\star\in K^V$ such that
\begin{equation}
(q_v^\star)^{-1}U_eq_u^\star=c_e
\quad\text{for all }e\in E^+,
\qquad h_C=\Id\quad\text{for all }C\in\mathcal C_T.
\label{eq:integrability}
\end{equation}
\end{definition}

Equation~\eqref{eq:integrability} is the finite common-clock condition.  It does
not identify node clock-orientation data with an externally imposed universal
time.  It asks only whether the matter-generated source transports and their
branch targets admit one compatible global section on the declared finite graph.

\section{One energy, two gravitational response branches}
\label{sec:energy}

\subsection{The finite synchronization energy}

For $k=(e^{i\theta},R)\in K$, define the phase character and full
clock-orientation discrepancy
\begin{equation}
\chi(k):=\cos\theta,
\qquad
\Psi(k):=1-\cos\theta+\frac{1}{2}\lVert R-I_3\rVert_F^2.
\label{eq:character}
\end{equation}
Then $\chi(\Id)=1$, $\chi(\zinv)=-1$,
$\chi(\zinv k)=-\chi(k)$, and $\Psi(k)\ge0$ with equality exactly at
$k=\Id$.  The character $\chi$ supplies the signed radial response term.  The
function $\Psi$ retains the full $\mathrm U(1)\times\mathrm{SO}(3)$ target,
transport, and cycle information in the alignment energy.

\begin{definition}[Finite clock-orientation energy]
\label{def:energy}
\Defn{} The complete energy is
\begin{align}
\Defect_{\mathbf p}(q,W,r)
={}&\sum_{e\in E^+}\left[
\frac{\alpha_e}{r_e^2}
-\frac{\gamma_e}{r_e}\chi\bigl(\mu_e(q,W)\bigr)
+\lambda_e\Psi\bigl(c_e^{-1}\mu_e(q,W)\bigr)
+\kappa_e\Psi\bigl(U_e^{-1}W_e\bigr)
\right] \nonumber\\
&+\sum_{C\in\mathcal C_T}\Lambda_C\Psi(h_C).
\label{eq:energy}
\end{align}
The final term fixes the declared cycle sector.  The first two terms give the
radial response; the third anchors the mismatch to the branch target selected by
Equation~\eqref{eq:pair-resolution}; and the fourth anchors actual transport to
matter-generated source transport.
\end{definition}

\begin{definition}[Complete flow and response]
\label{def:flow}
\Defn{} Equip the compact $K$ factors with a fixed product bi-invariant metric
and $(0,\infty)^{E^+}$ with its Euclidean metric.  The flow is
\begin{equation}
\dot q_v=-\grad_{q_v}\Defect_{\mathbf p},
\qquad
\dot W_e=-\grad_{W_e}\Defect_{\mathbf p},
\qquad
\dot r_e=-\partial_{r_e}\Defect_{\mathbf p}.
\label{eq:flow}
\end{equation}
Let $\Flow_{\mathbf p}(t,z_0)$ denote its solution.  Its radial response map is
\begin{equation}
\Resp_{\mathbf p}(q,W,r):=(\dot r_e)_{e\in E^+},
\qquad
\dot r_e=
\frac{2\alpha_e}{r_e^3}
-\frac{\gamma_e\chi(\mu_e)}{r_e^2}.
\label{eq:response}
\end{equation}
\end{definition}

\begin{proposition}[Finite dynamical closure]
\label{prop:closure}
\Thm{} For every admissible datum and every initial state in
$\State_{\mathbf p}$, the flow in Equation~\eqref{eq:flow} has a unique solution
for all $t\ge0$, remains in $\State_{\mathbf p}$, and satisfies
\begin{align}
\frac{d}{dt}\Defect_{\mathbf p}\bigl(\Flow_{\mathbf p}(t,z_0)\bigr)
={}&-\sum_{v\in V}\|\grad_{q_v}\Defect_{\mathbf p}\|^2
-\sum_{e\in E^+}\|\grad_{W_e}\Defect_{\mathbf p}\|^2
-\sum_{e\in E^+}|\partial_{r_e}\Defect_{\mathbf p}|^2
\le0.
\label{eq:descent}
\end{align}
\end{proposition}

\begin{proof}
The $K$ factors are compact and the energy is smooth.  For each edge,
$\alpha_e/r_e^2-\gamma_e\chi(\mu_e)/r_e\ge
\alpha_e/r_e^2-\gamma_e/r_e$, which tends to $+\infty$ as $r_e\downarrow0$
and is bounded below.  Energy descent therefore prevents a finite-time approach
to the excluded boundary $r_e=0$.  On every strip $r_e\ge\epsilon>0$, the
vector field is bounded on finite time intervals, so standard ODE extension gives
global existence and uniqueness.  Equation~\eqref{eq:descent} is the gradient-flow
identity.
\end{proof}

\subsection{Attractive, inverse, and frustrated sectors}

\begin{definition}[Physical branch predicates]
\label{def:branches}
\Defn{} For a source-target-integrable datum, define
\begin{align}
\Branch_+(\mathbf p):={}&\{z:\ h_C=\Id,\ c_e=\Id,\ W_e=U_e,
\ \mu_e=\Id\ \forall e,C\},
\label{eq:bplus}\\
\Branch_-(\mathbf p):={}&\{z:\ h_C=\Id,\ c_e=\zinv,\ W_e=U_e,
\ \mu_e=\zinv,\ r_e>\gamma_e/\lambda_e\ \forall e,C\},
\label{eq:bminus}\\
\Branch_{\mathrm{fr}}(\mathbf p):={}&\{z\in\State_{\mathbf p}:\exists C\in\mathcal C_T,
\ h_C\ne\Id\}.
\label{eq:bfr}
\end{align}
The labels are physical branch labels.  By Hypothesis~\ref{hyp:pair-resolution},
$\Branch_+$ is selected by physically co-registered counterposed pairs and
$\Branch_-$ by physically resolved counterposed pairs.
\end{definition}

\begin{proposition}[Forced attractive and inverse response]
\label{prop:response}
\Thm{} The branches in Definition~\ref{def:branches} are invariant under the
complete flow.  On the attractive branch,
\begin{equation}
\dot r_e=\frac{2\alpha_e-\gamma_er_e}{r_e^3}.
\label{eq:attractive-response}
\end{equation}
Thus $r_e^\star=2\alpha_e/\gamma_e$ is a stable radial equilibrium: $\dot r_e<0$
for $r_e>r_e^\star$ and $\dot r_e>0$ for $0<r_e<r_e^\star$.  On the inverse
branch,
\begin{equation}
\dot r_e=\frac{2\alpha_e+\gamma_er_e}{r_e^3}>0.
\label{eq:inverse-response}
\end{equation}
Moreover, the inverse mismatch target is a strict local minimum of the
edge energy on the mismatch quotient exactly when
\begin{equation}
r_e>\frac{\gamma_e}{\lambda_e}.
\label{eq:inverse-stability}
\end{equation}
\end{proposition}

\begin{proof}
On either aligned branch, the mismatch and transport penalty terms vanish with
their gradients.  In $\Branch_+$, $\chi(\mu_e)=1$; substituting into
Equation~\eqref{eq:response} gives Equation~\eqref{eq:attractive-response}.
In $\Branch_-$, $\chi(\mu_e)=\chi(\zinv)=-1$; substituting gives
Equation~\eqref{eq:inverse-response}.  The sign and stability statement for the
regular radial equilibrium follow by direct differentiation.

For the inverse target, write $\mu_e=\zinv\zeta$ near $\zeta=\Id$.  The
mismatch-dependent part of the edge energy is, up to terms fixed on the local
quotient,
\begin{equation}
-\frac{\gamma_e}{r_e}\chi(\mu_e)
+\lambda_e\Psi(c_e^{-1}\mu_e)
=
\lambda_e+
\left(\lambda_e-\frac{\gamma_e}{r_e}\right)\chi(\mu_e)
+\frac{\lambda_e}{2}\lVert R(\mu_e)-I_3\rVert_F^2.
\label{eq:inverse-local-energy}
\end{equation}
The phase coefficient in Equation~\eqref{eq:inverse-local-energy} is positive
exactly under Equation~\eqref{eq:inverse-stability}; the additional orientation
part of $\Psi$ is nonnegative and vanishes only at the aligned orientation.
Hence the full $K$-valued target is a strict local minimum on the mismatch
quotient exactly under Equation~\eqref{eq:inverse-stability}.  The condition is
preserved because Equation~\eqref{eq:inverse-response} makes $r_e$ increase.
The cycle sector is fixed datum-level boundary data, so the frustrated sector is
also invariant.
\end{proof}

Proposition~\ref{prop:response} is the required physical determination.  Once an
admissible finite datum, a pair-resolution branch, and the state are specified,
the theory fixes the sign of the edge response.  The inverse branch is therefore
not introduced by a later observer's label and not by a force-sign substitution in
an otherwise unchanged first-order equation.  It is a different, stable target
sector of the same finite gravitational dynamics.

\section{A(R)-space: registration, resolution, and blind spots}
\label{sec:registration}

\subsection{Typed registration bridge}

\begin{definition}[Registration bridge and observer quotient]
\label{def:bridge}
\Defn{} Fix a declared material observer/protocol $R$.  Its registration carrier
$\mathcal A_{\mathbf p,R}$ is a tagged record space.  A physical-to-registration
bridge is a total map
\begin{equation}
\Bridge_R:\State_{\mathbf p}\longrightarrow\mathcal A_{\mathbf p,R},
\qquad
z\longmapsto
\Reg_R(z):=\left(
\operatorname{tag}(\mathbf p),
\{\widehat\mu_{e,R},\widehat r_{e,R},\widehat\delta_{e,R}\}_{e\in E^+},
\widehat h_R,\mathcal N_R
\right),
\label{eq:bridge}
\end{equation}
where the hats denote registered quantities and $\mathcal N_R$ is the declared
instrument/noise and retention budget.  A quotient
\begin{equation}
\Quot_R:\mathcal A_{\mathbf p,R}\longrightarrow\overline{\ARspace}_{\mathbf p,R}
\label{eq:quotient}
\end{equation}
retains only the material record available under the protocol.  The composed
observation map is
\begin{equation}
\Obs_R:=\Quot_R\circ\Bridge_R.
\label{eq:observation-map}
\end{equation}
\end{definition}

\begin{definition}[Registration-side resolution]
\label{def:registration-resolution}
\Defn{} Let $\delta_R>0$ be the protocol's declared discrimination threshold and
$\epsilon_R\ge0$ a certified error bound for the registered pair-separation
coordinate.  The observer resolves the pair on edge $e$ at state $z$ only when
\begin{equation}
\Resolve_R(e;z)=1
\quad\Longleftrightarrow\quad
\widehat\delta_{e,R}(z)-\epsilon_R>\delta_R.
\label{eq:registration-resolution}
\end{equation}
Otherwise $\Resolve_R(e;z)=0$.  This is an A(R)-space condition.  It may differ
from the P-space condition $\delta_e>\eta_e$ in
Equation~\eqref{eq:pair-resolution}.
\end{definition}

There are therefore two thresholds with different types:
\begin{equation}
\underbrace{\eta_e}_{\text{intrinsic P-space branch threshold}}
\qquad\text{and}\qquad
\underbrace{(\delta_R,\epsilon_R)}_{\text{A(R)-space registration budget}}.
\label{eq:two-thresholds}
\end{equation}
A worse detector can alter $\Resolve_R$ but cannot alter $c_e$, the physical
branch target selected by Equation~\eqref{eq:pair-resolution}.

\subsection{The decision criterion}

\begin{definition}[Physical inverse target]
\label{def:target}
\Defn{} The physical inverse target is the branch predicate
\begin{equation}
\Theta_{\mathrm{inv}}(z)=
\begin{cases}
1,& z\in\Branch_-(\mathbf p),\\
0,& z\notin\Branch_-(\mathbf p).
\end{cases}
\label{eq:inverse-target}
\end{equation}
An observation map decides this target only when there is a function
$\overline\Theta_{\mathrm{inv},R}$ such that
\begin{equation}
\Theta_{\mathrm{inv}}
=
\overline\Theta_{\mathrm{inv},R}\circ\Obs_R.
\label{eq:factorization}
\end{equation}
\end{definition}

\begin{theorem}[Fiber criterion for the inverse branch]
\label{thm:fiber}
\Thm{} If there are states $z_+\in\Branch_+(\mathbf p)$ and
$z_-\in\Branch_-(\mathbf p)$ such that
\begin{equation}
\Obs_R(z_+)=\Obs_R(z_-),
\label{eq:co-registration}
\end{equation}
then no function $\overline\Theta_{\mathrm{inv},R}$ satisfying
Equation~\eqref{eq:factorization} exists.  Consequently, the material record
available through $\Obs_R$ cannot decide whether the physical state is in the
inverse branch.
\end{theorem}

\begin{proof}
If Equation~\eqref{eq:factorization} held, then equality of the two observations
would imply
$\Theta_{\mathrm{inv}}(z_+)=\Theta_{\mathrm{inv}}(z_-)$.
But Definition~\ref{def:target} gives
$\Theta_{\mathrm{inv}}(z_+)=0$ and
$\Theta_{\mathrm{inv}}(z_-)=1$, a contradiction.
\end{proof}

\begin{corollary}[Resolved P-space branch, unresolved material record]
\label{cor:blindspot}
\Cnd{} Suppose $z_-\in\Branch_-(\mathbf p)$, so that the pair is physically
resolved in P-space, and suppose its material record satisfies
$\Resolve_R(e;z_-)=0$.  If the quotient identifies this record with an
attractive-branch record as in Equation~\eqref{eq:co-registration}, then the
inverse branch is physically present but unresolvable by that material
observation map.
\end{corollary}

\begin{corollary}[Energy scaling inside an unchanged quotient]
\label{cor:energy}
\Cnd{} Let $E$ label a family of protocols whose retained quotient is invariant:
$\Obs_{R,E}=\Obs_R$ for all declared $E$.  If the hypothesis of
Theorem~\ref{thm:fiber} holds for $\Obs_R$, then increasing $E$ within that
unchanged observation quotient cannot make $\Theta_{\mathrm{inv}}$ factor
through the record.  A changed energy scale helps only if it changes the bridge
or quotient so that the relevant fiber is separated.
\end{corollary}

\section{The second-order theory object}
\label{sec:secondorder}

\begin{definition}[Finite second-order inverse-gravity theory]
\label{def:secondorder}
\Defn{} For an admissible finite datum $\mathbf p$ and a declared registration
protocol $R$, the second-order theory object is
\begin{equation}
\mathfrak G^{(2)}_{\mathbf p,R}:=
\left(
\underbrace{\Pspace_{\mathbf p},\State_{\mathbf p},\Defect_{\mathbf p},
\Flow_{\mathbf p},\Resp_{\mathbf p},\Branch_+,\Branch_-,\Branch_{\mathrm{fr}}}_{\text{first-order gravitational carrier}},
\underbrace{\mathcal A_{\mathbf p,R},\Bridge_R,\Quot_R,\Obs_R,\Resolve_R}_{\text{registration theory}},
\Theta_{\mathrm{inv}}
\right).
\label{eq:second-order-object}
\end{equation}
It is a bridged composition, not an identification, of a first-order gravitational
carrier in P-space and a distinct theory of registration in $A(R)$-space.
\end{definition}

\begin{theorem}[Second-order inverse-gravity result]
\label{thm:secondorder}
\Thm{} For every admissible finite datum $\mathbf p$ and declared total bridge
$\Bridge_R$, the object in Equation~\eqref{eq:second-order-object} has the
following properties.
\begin{enumerate}
\item The P-space carrier determines a regular attractive response on
$\Branch_+$ and a positive inverse response on $\Branch_-$, by
Proposition~\ref{prop:response}.
\item The branch condition is physically typed by the intrinsic pair-resolution
law, not by the observer's detector budget.
\item The A(R)-space component decides the inverse target exactly when the
factorization condition in Equation~\eqref{eq:factorization} holds.
\item Where the fiber condition of Theorem~\ref{thm:fiber} holds, the theory
proves a target-relative material observability blind spot without treating the
physical branch as absent.
\end{enumerate}
\end{theorem}

\begin{proof}
Items 1 and 2 are respectively Proposition~\ref{prop:response} and
Hypothesis~\ref{hyp:pair-resolution}.  Item 3 is Definition~\ref{def:target}.
Item 4 is Theorem~\ref{thm:fiber}.  Each inference uses a named object within
one space or an explicit bridge between the spaces.
\end{proof}


% =================================================================
% Operand-level fitness self-selection insertion (v1)
% =================================================================

% ============================================================================
% INSERT FOR SUPPLEMENTARY MATERIALS
% Suggested placement: after Section~\ref{sec:secondorder} and before the
% implementation-schedule section.  This insert assumes the theorem
% environments and the macros \Pspace, \ARspace, \Bridge, \Status, \Defn,
% \Hyp, \Thm, and \Cnd already defined in the supplement preamble.
% ============================================================================

% ============================================================================
% INSERT FOR SUPPLEMENTARY MATERIALS (v5)
% Replace the prior operand-level self-selection insert with this section.
% Suggested placement: after Section~\ref{sec:secondorder} and before
% \section{Schedules are implementation conditions}.
%
% This insert assumes the supplement preamble already defines \Pspace,
% \ARspace, \Bridge, \Defn, \Hyp, \Thm, \Cnd, and theorem environments.
% ============================================================================

\section{Operand-level viability closure and physical self-selection}
\label{sec:operand-self-selection-v5}

First-order physics normally seeks a law that explains the matter we observe.
The proposed inversion in second-order physics is that energy functionals and
matter mutually select one another: a fluctuation can become physically
consequential only when what it produces creates the conditions required for
that same physical process to continue.

This section states that inversion in a typed form.  It does not treat fitness
as an additive physical-energy contribution, and it does not make a claim about
theories proposed by human observers.  The relevant candidate class is instead
the declared class of fluctuation-induced physical or registration-capable
executions.  The finite fixed-point construction is performed on a lattice of
fitness/viability certificates, not on the continuous P-space carrier.  It therefore
proves certificate closure once a candidate execution is supplied; the claim
that physical operativity itself obeys this closure is stated separately as a
physical hypothesis.

\begin{definition}[Operative fluctuation candidate and realization map]
\label{def:ossv5-operative-candidate}
\Defn{} Fix a target $T$ and resolution $\delta$.  Let
\begin{equation}
\mathcal Q_{\mathrm{op};T,\delta}
\label{eq:ossv5-qop}
\end{equation}
be a declared class of fluctuation-induced candidates capable, in the declared
model domain, of producing a matter-bearing P-space continuation, an $A(R)$-space
registration, or both.  This is not a space of human-authored theories.

A realization map is a typed, model-dependent relation
\begin{equation}
\mathfrak I_{T,\delta}:
\mathcal Q_{\mathrm{op};T,\delta}
\longrightarrow
\bigsqcup_{h\in\mathcal H_{T,\delta}}
\bigl(\{h\}\times\mathcal X_{T,\delta}\bigr),
\qquad
q\longmapsto(h_q,x_q),
\label{eq:ossv5-realization-map}
\end{equation}
where $h_q$ is the candidate energy-functional execution induced by $q$ and
$x_q$ is its typed matter--registration base record.  Equation~
\eqref{eq:ossv5-realization-map} does not identify a fluctuation with a
Hamiltonian; it records the candidate functional execution and operand induced
by that fluctuation in the declared model.
\end{definition}

\begin{hypothesis}[Physical operativity through endogenous closure]
\label{hyp:ossv5-physical-operativity}
\Hyp{} A candidate $q\in\mathcal Q_{\mathrm{op};T,\delta}$ becomes physically
or registrationally consequential only if, for
$\mathfrak I_{T,\delta}(q)=(h_q,x_q)$, the candidate execution is admitted on
its closed viability-bearing operand:
\begin{equation}
\mathsf{Adm}_{h_q,T,\delta}\!\left(\Xi_{h_q}^{\star}(x_q)\right)=1.
\label{eq:ossv5-operativity-condition}
\end{equation}
This is an ontic hypothesis of the proposed second-order physical theory.  It
is not a consequence of the fixed-point theorem below, and it does not assert
that every fluctuation is realized or that the admitted functional class is
universally unique.
\end{hypothesis}

\begin{definition}[Typed base record and viability-bearing operand]
\label{def:ossv5-operand}
\Defn{} Let
\begin{equation}
\mathcal X_{T,\delta}
:=
\Pspace_{T,\delta}
\times
\left(\ARspace\right)_{T,\delta}
\times
\mathcal B_{T,\delta}
\times
\mathcal K_{T,\delta}
\label{eq:ossv5-base-record}
\end{equation}
be a typed record product.  Its coordinates respectively retain a physical
state/path, a registration state/path, a declared bridge record, and a
history--regime--support record.  Equation~\eqref{eq:ossv5-base-record} is a
record constructor; it is not an identification of the spaces named by its
coordinates.

Let $\mathcal V_{T,\delta}$ be a finite complete lattice of fitness/viability
certificates, ordered by declared certificate inclusion $\preceq$, with least
element $\bot$.  A certificate $\phi\in\mathcal V_{T,\delta}$ may retain
only the finite predicates required by the declared model class---for example,
continuation, balance, support, bridge definition, regime preservation,
distinguishability, recovery, and retention.  It is not a physical field, a
force, an additive term in the energy, or an external selector.

For a candidate energy functional $h\in\mathcal H_{T,\delta}$, a
viability-bearing operand is
\begin{equation}
\Xi^h_{T,\delta}=(x,\phi)
\in
\mathcal X_{T,\delta}\times\mathcal V_{T,\delta}.
\label{eq:ossv5-operand}
\end{equation}
\end{definition}

\begin{hypothesis}[Endogenous certificate update]
\label{hyp:ossv5-certificate-update}
\Hyp{} For every $h\in\mathcal H_{T,\delta}$ and every declared base record
$x\in\mathcal X_{T,\delta}$, the coupled P-space update, $A(R)$-space update,
and typed bridge consequences of $h$ induce a map
\begin{equation}
\Phi_{h,x}:\mathcal V_{T,\delta}\longrightarrow\mathcal V_{T,\delta}.
\label{eq:ossv5-certificate-update}
\end{equation}
The map is monotone:
\begin{equation}
\phi\preceq\psi
\quad\Longrightarrow\quad
\Phi_{h,x}(\phi)\preceq\Phi_{h,x}(\psi).
\label{eq:ossv5-monotonicity}
\end{equation}
The map records viability consequences generated under $h$ itself.  It
introduces no fitness judge external to the candidate functional and its typed
physical--registration execution.
\end{hypothesis}

\begin{definition}[Closed certificate, full admissibility, and surviving class]
\label{def:ossv5-admissibility}
\Defn{} Under Hypothesis~\ref{hyp:ossv5-certificate-update}, define the
Kleene certificate sequence
\begin{equation}
\phi^{(0)}_{h,x}:=\bot,
\qquad
\phi^{(n+1)}_{h,x}:=\Phi_{h,x}\!\left(\phi^{(n)}_{h,x}\right).
\label{eq:ossv5-kleene-sequence}
\end{equation}
Whenever this sequence stabilizes, write its stable value as
$\phi_h^\star(x)$ and define the closed operand
\begin{equation}
\Xi_h^\star(x):=\bigl(x,\phi_h^\star(x)\bigr).
\label{eq:ossv5-closed-operand}
\end{equation}

For each $h$, let
\begin{equation}
\mathsf{Adm}^{\mathrm{viab}}_{h,T,\delta}:
\mathcal X_{T,\delta}\times\mathcal V_{T,\delta}
\longrightarrow\{0,1\}
\label{eq:ossv5-viability-admissibility}
\end{equation}
be the viability component of admissibility.  Its arguments include the
certificate generated under that same candidate $h$.

For a gravitational candidate $h$ and an edge $e$ on its regular aligned
branch, write the radial response as
\begin{equation}
f^+_{h,e}(r):=
\frac{2\alpha_e^h-\gamma_e^h r}{r^3},
\qquad r>0.
\label{eq:ossv5-regular-response}
\end{equation}
The carrier-regularity predicate is
\begin{equation}
\mathsf{Reg}_{h,T,\delta}(x)=1
\quad\Longleftrightarrow\quad
\forall e\in E^+\;\exists r^\star_{h,e}>0:\
 f^+_{h,e}(r^\star_{h,e})=0
\ \wedge\
 \bigl(f^+_{h,e}\bigr)'(r^\star_{h,e})<0.
\label{eq:ossv5-carrier-regularity}
\end{equation}
Thus regularity requires a positive locally stable equilibrium on the regular
branch.  It contains no inverse-branch response-sign condition.

The full admissibility predicate is the conjunction
\begin{equation}
\mathsf{Adm}^{\mathrm{full}}_{h,T,\delta}(x,\phi)=1
\quad\Longleftrightarrow\quad
\mathsf{Reg}_{h,T,\delta}(x)=1
\ \wedge\
\mathsf{Adm}^{\mathrm{viab}}_{h,T,\delta}(x,\phi)=1.
\label{eq:ossv5-full-admissibility}
\end{equation}
An evaluation is admitted precisely when
\begin{equation}
H_h\!\left[\Xi_h^\star(x)\right]
\quad\text{is evaluated only if}\quad
\mathsf{Adm}^{\mathrm{full}}_{h,T,\delta}\!\left(\Xi_h^\star(x)\right)=1.
\label{eq:ossv5-admitted-evaluation}
\end{equation}

The surviving functional class at $x$ is
\begin{equation}
\mathcal H_{\mathrm{adm}}(x):=
\left\{
 h\in\mathcal H_{T,\delta}:
 \mathsf{Adm}^{\mathrm{full}}_{h,T,\delta}\!\left(\Xi_h^\star(x)\right)=1
\right\}.
\label{eq:ossv5-surviving-class}
\end{equation}
\end{definition}

\begin{theorem}[Finite endogenous certificate closure]
\label{thm:ossv5-finite-closure}
\Thm{} Under Hypothesis~\ref{hyp:ossv5-certificate-update}, for every
$h\in\mathcal H_{T,\delta}$ and $x\in\mathcal X_{T,\delta}$, the sequence in
Equation~\eqref{eq:ossv5-kleene-sequence} is ascending and stabilizes after at
most $|\mathcal V_{T,\delta}|-1$ strict increases.  Its stable value
$\phi_h^\star(x)$ is the least fixed point of $\Phi_{h,x}$.  Therefore the
closed operand in Equation~\eqref{eq:ossv5-closed-operand} exists for every
candidate and base record on the declared finite certificate domain.
\end{theorem}

\begin{proof}
Because $\bot$ is the least element of $\mathcal V_{T,\delta}$,
$\phi^{(0)}_{h,x}\preceq\phi^{(1)}_{h,x}$.  If
$\phi^{(n)}_{h,x}\preceq\phi^{(n+1)}_{h,x}$, monotonicity
(Equation~\eqref{eq:ossv5-monotonicity}) gives
\begin{equation}
\phi^{(n+1)}_{h,x}
=
\Phi_{h,x}\!\left(\phi^{(n)}_{h,x}\right)
\preceq
\Phi_{h,x}\!\left(\phi^{(n+1)}_{h,x}\right)
=
\phi^{(n+2)}_{h,x}.
\end{equation}
Induction proves that the sequence is ascending.  A strictly ascending chain in
the finite set $\mathcal V_{T,\delta}$ has at most
$|\mathcal V_{T,\delta}|-1$ strict increases.  Hence there is an index $N$ with
$\phi^{(N)}_{h,x}=\phi^{(N+1)}_{h,x}$, so
$\phi^{(N)}_{h,x}$ is a fixed point.

Let $\psi$ be any fixed point of $\Phi_{h,x}$.  Since $\bot\preceq\psi$,
monotonicity yields, by induction,
$\phi^{(n)}_{h,x}\preceq\psi$ for all $n$.  In particular,
$\phi^{(N)}_{h,x}\preceq\psi$.  Thus the stable value is below every fixed
point and is the least fixed point.  Equation~\eqref{eq:ossv5-closed-operand}
then defines a closed operand for every $h$ and $x$.
\end{proof}

\begin{theorem}[Endogenous admissibility is operand-level, not an energy term]
\label{thm:ossv5-endogenous-admissibility}
\Thm{} For every $h$ and $x$, membership
$h\in\mathcal H_{\mathrm{adm}}(x)$ holds if and only if the candidate's own
coupled consequences generate the certificate $\phi_h^\star(x)$ and its
admissibility predicate accepts the resulting operand.  Consequently, the
construction retains a candidate functional only through the viability-bearing
operand generated under that candidate: it does not introduce a fitness-energy
term or an external functional selector.  It returns a unique selected
functional exactly when
\begin{equation}
\left|\mathcal H_{\mathrm{adm}}(x)\right|=1;
\label{eq:ossv5-uniqueness}
\end{equation}
otherwise it returns the surviving class.
\end{theorem}

\begin{proof}
By Definition~\ref{def:ossv5-admissibility}, the certificate
$\phi_h^\star(x)$ is generated by iterating the $h$-induced map
$\Phi_{h,x}$, and membership in Equation~\eqref{eq:ossv5-surviving-class} is
exactly the truth of
$\mathsf{Adm}^{\mathrm{full}}_{h,T,\delta}(x,\phi_h^\star(x))=1$.
The viability component is evaluated on the certificate generated by $h$'s own
coupled consequences, while the regularity component is evaluated on $h$'s own
regular-branch response.  Neither operation adds a fitness term to $H_h$ or
introduces a selector external to the candidate functional and its typed
execution.  The last statement follows directly from the cardinality of the
surviving class.
\end{proof}

\begin{theorem}[Regular-branch admissibility forces coefficient positivity]
\label{lem:ossv5-regularity-positivity}
\Thm{} Let $h\in\mathcal H_{\mathrm{adm}}(x)$ and let $e\in E^+$.
Then the carrier coefficients on the corresponding regular branch satisfy
\begin{equation}
\alpha_e^h>0,
\qquad
\gamma_e^h>0.
\label{eq:ossv5-coefficient-positivity}
\end{equation}
The conclusion follows from the carrier-regularity component of full
admissibility; it does not use an inverse-branch response sign.
\end{theorem}

\begin{proof}
Since $h\in\mathcal H_{\mathrm{adm}}(x)$,
Equation~\eqref{eq:ossv5-full-admissibility} gives
$\mathsf{Reg}_{h,T,\delta}(x)=1$.  Thus there is an
$r^\star_{h,e}>0$ such that
\begin{equation}
0=f^+_{h,e}(r^\star_{h,e})
=
\frac{2\alpha_e^h-\gamma_e^h r^\star_{h,e}}
{(r^\star_{h,e})^3},
\end{equation}
so
\begin{equation}
\gamma_e^h=\frac{2\alpha_e^h}{r^\star_{h,e}}.
\label{eq:ossv5-gamma-from-alpha}
\end{equation}
Differentiating Equation~\eqref{eq:ossv5-regular-response} gives
\begin{equation}
\bigl(f^+_{h,e}\bigr)'(r)
=-\frac{6\alpha_e^h}{r^4}+\frac{2\gamma_e^h}{r^3}.
\end{equation}
Substituting Equation~\eqref{eq:ossv5-gamma-from-alpha} at
$r=r^\star_{h,e}$ yields
\begin{equation}
\bigl(f^+_{h,e}\bigr)'(r^\star_{h,e})
=-\frac{2\alpha_e^h}{(r^\star_{h,e})^4}.
\end{equation}
Carrier regularity makes the left-hand side strictly negative.  Because
$r^\star_{h,e}>0$, this implies $\alpha_e^h>0$.  Equation~\eqref{eq:ossv5-gamma-from-alpha}
then gives $\gamma_e^h>0$.
\end{proof}

\begin{remark}[Logical role of the regularity condition]
\label{rem:ossv5-regularity-role}
\Cnd{} In the finite carrier developed earlier in this supplement, positivity
of the coefficients is part of the declared carrier data.  The present
formulation makes explicit a stronger candidate-class criterion: when full
admissibility requires a positive locally stable regular-branch equilibrium,
coefficient positivity follows algebraically from that criterion.  This is a
reorganization of the candidate-class assumption, not an empirical derivation
of coefficient signs from the inverse branch.
\end{remark}

\begin{definition}[Viability-distinguishing fiber]
\label{def:ossv5-fitness-fiber}
\Defn{} Let
\begin{equation}
\pi_{\mathcal X}:
\mathcal X_{T,\delta}\times\mathcal V_{T,\delta}
\longrightarrow
\mathcal X_{T,\delta},
\qquad
\pi_{\mathcal X}(x,\phi)=x
\label{eq:ossv5-fitness-erasure}
\end{equation}
be the certificate-erasing projection.  The admissibility family is
\emph{viability-sensitive} when there are $h$, $x$, and
$\phi_0,\phi_1\in\mathcal V_{T,\delta}$ such that
\begin{equation}
\mathsf{Adm}^{\mathrm{full}}_{h,T,\delta}(x,\phi_0)
\neq
\mathsf{Adm}^{\mathrm{full}}_{h,T,\delta}(x,\phi_1).
\label{eq:ossv5-viability-fiber}
\end{equation}
\end{definition}

\begin{theorem}[Necessity of viability-bearing information]
\label{thm:ossv5-nonreducibility}
\Thm{} If the viability-distinguishing fiber condition in
Equation~\eqref{eq:ossv5-viability-fiber} holds, then
$\mathsf{Adm}^{\mathrm{full}}_{h,T,\delta}$ does not factor through
$\pi_{\mathcal X}$.  Hence no exact endogenous admissibility decision retaining
only the fitness-erased base record can implement the declared viability-sensitive
selection.

More generally, a retained representation
$\rho:\mathcal X_{T,\delta}\times\mathcal V_{T,\delta}\to Y$ supports exact
admissibility selection if and only if, for every $h$, there is a map
$d_h:Y\to\{0,1\}$ satisfying
\begin{equation}
\mathsf{Adm}^{\mathrm{full}}_{h,T,\delta}=d_h\circ\rho
\quad\text{on the declared admissibility domain.}
\label{eq:ossv5-exact-factorization}
\end{equation}
Thus an explicit certificate coordinate may be replaced only by a faithful
sufficient statistic for every admissibility-relevant viability distinction.
\end{theorem}

\begin{proof}
Suppose, to the contrary, that
$\mathsf{Adm}^{\mathrm{full}}_{h,T,\delta}=\bar d_h\circ\pi_{\mathcal X}$ for some
$\bar d_h:\mathcal X_{T,\delta}\to\{0,1\}$.  Applying this equality to the
two operands in Equation~\eqref{eq:ossv5-viability-fiber} gives
\begin{equation}
\mathsf{Adm}^{\mathrm{full}}_{h,T,\delta}(x,\phi_0)
=
\bar d_h(x)
=
\mathsf{Adm}^{\mathrm{full}}_{h,T,\delta}(x,\phi_1),
\end{equation}
contradicting Equation~\eqref{eq:ossv5-viability-fiber}.  Therefore no such
factorization exists.

For the general statement, if Equation~\eqref{eq:ossv5-exact-factorization}
holds, evaluation of $d_h(\rho(x,\phi))$ reproduces the declared admissibility
value exactly.  Conversely, suppose that $\rho$ supports exact selection.  For
each $y$ in the image of the declared domain, choose any operand
$(x,\phi)$ with $\rho(x,\phi)=y$ and define
$d_h(y):=\mathsf{Adm}^{\mathrm{full}}_{h,T,\delta}(x,\phi)$.  Exactness ensures that this is
independent of the chosen representative: otherwise one $\rho$-fiber would
contain both an admitted and a rejected operand and no decision based only on
$y$ could be exact.  Hence $d_h$ is well defined and satisfies
Equation~\eqref{eq:ossv5-exact-factorization}.
\end{proof}

\begin{hypothesis}[Inverse branch and sign-blind viability admission]
\label{hyp:ossv5-sign-neutrality}
\Hyp{} For every $h\in\mathcal H_{\mathrm{adm}}(x)$ on an inverse branch,
the carrier branch relations give
\begin{equation}
\dot r_e^h=
\frac{2\alpha_e^h+\gamma_e^h r_e}{r_e^3},
\qquad r_e>0.
\label{eq:ossv5-inverse-response}
\end{equation}
The viability component
$\mathsf{Adm}^{\mathrm{viab}}_{h,T,\delta}$ contains no inverse-response-sign
predicate and does not retain or reject a candidate by requiring
$\mathsf{sgn}(\dot r_e^h)=+1$.  By Definition~\ref{def:ossv5-admissibility},
the regularity component tests only the regular branch in
Equation~\eqref{eq:ossv5-regular-response}; it likewise contains no
inverse-response-sign criterion.
\end{hypothesis}

\begin{theorem}[Sign invariance over the surviving functional class]
\label{thm:ossv5-sign-invariance}
\Thm{} Under Hypothesis~\ref{hyp:ossv5-sign-neutrality}, every admitted
candidate has the same positive inverse-branch response sign:
\begin{equation}
\forall h\in\mathcal H_{\mathrm{adm}}(x),
\qquad
\mathsf{sgn}(\dot r_e^h)=+1.
\label{eq:ossv5-sign-invariance}
\end{equation}
Thus the response sign is invariant over the surviving class.  Its positivity
is derived from the carrier-regularity component of full admissibility, not
stipulated as an inverse-sign selection rule.
\end{theorem}

\begin{proof}
Let $h\in\mathcal H_{\mathrm{adm}}(x)$.  By
Lemma~\ref{lem:ossv5-regularity-positivity},
$\alpha_e^h>0$ and $\gamma_e^h>0$.  Hypothesis~\ref{hyp:ossv5-sign-neutrality}
gives $r_e>0$.  Hence the denominator in
Equation~\eqref{eq:ossv5-inverse-response} is positive and its numerator
$2\alpha_e^h+\gamma_e^h r_e$ is strictly positive.  Therefore
$\dot r_e^h>0$, proving Equation~\eqref{eq:ossv5-sign-invariance}.  Since the
argument applies to every admitted $h$, the sign is constant on the surviving
class.

Neither component of full admissibility contains the conclusion
$\mathsf{sgn}(\dot r_e^h)=+1$: the viability component is sign-blind by
Hypothesis~\ref{hyp:ossv5-sign-neutrality}, and the regularity component is
stated entirely on the regular branch.  The inverse response sign is therefore
an algebraic consequence after admission, not a premise used to obtain
admission.  Hypothesis~\ref{hyp:ossv5-physical-operativity} is not used in this
proof.
\end{proof}

\begin{remark}[Status and scope]
\label{rem:ossv5-status}
\Cnd{} The formal results in this section establish finite certificate closure,
operand-level endogenous admissibility, the necessity of preserving
admissibility-relevant viability information, and sign invariance over the
surviving class under the stated carrier and branch hypotheses.  In particular,
full admissibility requires a positive locally stable regular-branch equilibrium;
Lemma~\ref{lem:ossv5-regularity-positivity} then derives the coefficient
positivity used by the inverse-response theorem.  The derivation is
model-internal and does not use the physical-operativity hypothesis.

The physical-operativity hypothesis is distinct.  It proposes that the relevant
candidate class consists of fluctuation-induced physical or registration-capable
executions, and that only those whose own products preserve the conditions
required for continuation become consequential.  That is an ontic hypothesis of
the proposed second-order theory, not a consequence of certificate closure and
not a claim that all human-proposed theories converge to one formula.  The
construction also does not prove that every fluctuation is realized or that the
surviving class is universally unique.
\end{remark}

\section{Schedules are implementation conditions}
\label{sec:schedules}

\begin{definition}[Target-conditioned inverse schedule]
\label{def:schedule}
\Defn{} Fix an inverse target state $x_\star\in\Branch_-(\mathbf p)$, an
admissible operation family $\mathcal O_P$, and a bridge/registration discipline
$B$.  A target-conditioned inverse schedule is a finite ordered record
\begin{equation}
\Sched_{\mathcal E,B}(x_\star)=
\left\langle
(u_k,\Delta t_k,\mathcal B_k,\mathcal R_k,S_k,B_k)
\right\rangle_{k=1}^n,
\label{eq:schedule}
\end{equation}
where $u_k\in\mathcal O_P$ is an ordinary forward-time admissible physical
operation, $\Delta t_k$ is its duration, $\mathcal B_k$ is its active boundary or
interface condition, $\mathcal R_k$ is its regime record, $S_k$ is its support,
and $B_k$ is its P-space-to-registration bridge.  The schedule is admissible only
when every transition composes, stays inside the declared carrier/regime corridor,
and terminates at or within a stated approximation of $x_\star$.
\end{definition}

\begin{proposition}[Schedule status]
\label{prop:schedule}
\Cnd{} A schedule in Definition~\ref{def:schedule} is a realization witness only
when it exists and satisfies its declared corridor.  It neither defines the
inverse branch nor changes its physical membership predicate.  Inverse membership
is fixed first by Equations~\eqref{eq:pair-resolution} and~\eqref{eq:bminus}; the
schedule is a separate conditional object for reaching a state in that branch.
\end{proposition}

\section{Observable consequences and falsification route}
\label{sec:falsification}

The theory already contains a configuration-level determination.  For every edge
on a declared branch, the response sign is fixed:
\begin{equation}
\mathsf{sgn}(\dot r_e)=
\begin{cases}
\mathsf{sgn}(2\alpha_e-\gamma_er_e),&z\in\Branch_+(\mathbf p),\\
+1,&z\in\Branch_-(\mathbf p).
\end{cases}
\label{eq:sign-prediction}
\end{equation}
This is not yet a claim about a continuum metric, astronomical body, or laboratory
apparatus.  It is a prediction within the finite carrier.  Any continuum or
experimental claim must supply a reconstruction map and an observable bridge.

\begin{definition}[Effective-regime reconstruction obligation]
\label{def:reconstruction}
\Open{} A continuum or other effective rendering requires a declared target
state space $X_\delta$, a dynamics $U_\delta(t)$ on that space, and a map
\begin{equation}
\mathsf S_{\mathbf p,\delta}:\State_{\mathbf p}\longrightarrow X_\delta
\label{eq:reconstruction}
\end{equation}
with a stated norm, metric, probability law, or retained quotient for which
\begin{equation}
\mathsf S_{\mathbf p,\delta}\circ\Flow_{\mathbf p,t}
\approx_{\varepsilon_{\mathrm{dyn}}(\delta)}
U_\delta(t)\circ\mathsf S_{\mathbf p,\delta}
\label{eq:effective-recovery}
\end{equation}
holds on a declared domain and interval.  A physical detector claim additionally
requires a bridge from the reconstructed observable to a registration record.
\end{definition}

\begin{obligation}[Branch-sensitive falsification]
\label{obl:falsification}
\Open{} A candidate physical realization of the finite carrier is falsified on a
declared inverse configuration if it simultaneously establishes: (i) the
P-space branch conditions in Equation~\eqref{eq:bminus}; (ii) the applicable
stability condition in Equation~\eqref{eq:inverse-stability}; and (iii) a valid
reconstruction/observable bridge for which the corresponding response has
$\dot r_e\le0$.  Conversely, a claimed observation of inverse gravity is not
licensed merely by a registration label: the branch and bridge records must be
shown.
\end{obligation}

\section{Claim boundary and conclusion}
\label{sec:conclusion}

The supplement establishes the following finite theory statement:
\begin{quote}
\emph{At a declared primitive scale, a counterposed gravitational interaction
pair is physically co-registered below an intrinsic synchronization threshold and
physically resolved above it.  The co-registered and resolved conditions select,
respectively, regular attractive and inverse-response sectors of one finite
clock-orientation gravity dynamics.  The latter has a forced positive separation
response under its stated stability condition.  A distinct P-space-to-A(R)-space
bridge determines whether a material record retains the branch distinction; when
the observation quotient identifies an attractive state with an inverse state,
the inverse target cannot be decided from that record.}
\end{quote}

The finite P-space carrier, its energy, its branch predicates, and its response
law are fully specified in this supplement.  The P-space/A(R)-space bridge and
factorization criterion make the second-order content explicit.  Physical
realization, continuum recovery, and empirical selection are not imported into
the theorem; they are the next typed obligations of the theory object.

\end{document}
